%------------------------------------------------------------------------------%
\documentclass[fleqn,utf8,aspectratio=169,12pt]{beamer}

\usepackage{integracao_e_series}
\usepackage{cancel}

\title{Séries Numéricas -- Definição}

%------------------------------------------------------------------------------%
\begin{document}

\addtitle

%------------------------------------------------------------------------------%
%------------------------------------------------------------------------------%
\section{Definição}

%------------------------------------------------------------------------------%
\begin{frame}{Séries}
    Uma \emph{Série Infinita} é a soma dos termos de uma \emph{Sequência Infinita} $(a_n)$

    \pause
    \vspace{\baselineskip}
    Se existir, essa \emph{soma} pode ser escrita como
    \[
      \sum_{n=1}^\infty a_n
      \;=\; a_1 + a_2 + \cdots + a_n + \cdots
      \;=\; S
    \]

    \pause
    Cada valor $a_n$ é chamado de \emph{Termo da Série}

		\only<beamer>{
    \pause
    \vspace{\baselineskip}
    \begin{center}
        \textbf{\color{red}Como saber se a soma existe?}
    \end{center}}
\end{frame}

%------------------------------------------------------------------------------%
\begin{frame}{Sequência de Somas Parciais}
    Soma dos primeiros $n$ \emph{termos} de uma sequência $(a_k)$
    \[
      S_{\memph{n}}
      = \; a_1 + a_2 + \cdots + a_{\memph{n}} \;
      = \sum_{k=1}^{\memph{n}} a_k
    \]

    \pause
    Temos agora uma nova sequência: \emph{Somas Parciais} $(S_n)$

    \pause
    \vspace{\baselineskip}
    Podemos usar as técnicas definidas para sequências

\end{frame}

%------------------------------------------------------------------------------%
\begin{frame}{Convergência de uma Série}
    A \emph{soma da série} é $S$
    \[
      \sum_{n=1}^\infty a_n = S
    \]
    quando a sequência de \emph{somas parciais} $(S_n)$ converge para $S$
    \pause
    \[
      \lim_{n\to\infty} S_n
      \;=\; \lim_{n\to\infty} \sum_{k=1}^n a_k
      \;=\; S
    \]
    \pause
    Se sequência de somas parciais diverge a série \emph{diverge}
\end{frame}

%------------------------------------------------------------------------------%
\begin{frame}{Exemplo}
    Somar os termos da sequência
    \[ a_k = \frac{1}{2^k} \]

    \pause
    \vspace{.5\baselineskip}
    \[
    	a_1 = \frac{1}{2}  \qquad\pause
    	a_2 = \frac{1}{4}  \qquad\pause
    	a_3 = \frac{1}{8}  \qquad\pause
      a_4 = \frac{1}{16} \qquad\pause
      a_5 = \frac{1}{32} \qquad\pause\cdots
    \]
\end{frame}

%------------------------------------------------------------------------------%
\begin{frame}{Somas Parciais}
    \begin{align*}
      \onslide<+->{a_1 & = \frac{1}{2}   & }
      \onslide<+->{S_1 & = \frac{1}{2}                             &  & }
      \onslide<+->{\hspace{-17mm}  = 1 - \frac{1}{2}   \\[2mm]}
    	\onslide<+->{a_2 & = \frac{1}{4}   & }
      \onslide<+->{S_2 & = \frac{1}{2} + \frac{1}{4}               &  & }
      \onslide<+->{\hspace{-17mm}  = 1 - \frac{1}{4}   \\[2mm]}
    	\onslide<+->{a_3 & = \frac{1}{8}   & }
      \onslide<+->{S_3 & = \frac{1}{2} + \frac{1}{4} + \frac{1}{8} &  & }
      \onslide<+->{\hspace{-17mm}  = 1 - \frac{1}{8}   \\[2mm]}
    	\onslide<+->{a_4 & = \frac{1}{16}  & }
      \onslide<+->{S_4 & = \frac{1}{2} + \frac{1}{4} + \frac{1}{8} + \frac{1}{16} &  & }
      \onslide<+->{\hspace{-17mm}  = 1 - \frac{1}{16}   \\[7mm]}
      \onslide<+->{a_n & = \frac{1}{2^n} & }
      \onslide<+->{S_n & = \sum_{k=1}^{n} \frac{1}{2^k}        &  & }
      \onslide<+->{\hspace{-17mm}  = \emph{1 - \frac{1}{2^n} \quad \text{(Conjectura)} }}
    \end{align*}
\end{frame}

%------------------------------------------------------------------------------%
\begin{frame}{Demonstração dor Indução Finita}
  Demonstração dor Indução Finita

  \pause
  \vspace{\baselineskip}
	Para mostrar que uma afirmação é verdadeira para todo $n\in\N$

  devemos provar que
  \pause
	\begin{enumerate}[<+->]
		\item ela é verdadeira para $n=1$
		\item se ela for verdadeira para $n$ então também será para $n+1$
	\end{enumerate}
\end{frame}

%------------------------------------------------------------------------------%
\begin{frame}{Demonstração da Fórmula das Somas Parciais}
\begin{columns}
\column{0.5\linewidth}
  Assumimos por \emph{indução finita} que
  \quad
  \[
    S_{n-1} = 1 - \frac{1}{2^{n-1}}
  \]
  para um \emph{$n$ dado qualquer}

  \pause
  \vspace{\baselineskip}
  Queremos mostrar que
  \[
    S_n = 1 - \frac{1}{2^n}
  \]

  \pause
  \column{0.5\linewidth}
  \begin{align*}
    \onslide<+->{S_n & = \frac{1}{2}
                       + \frac{1}{4}
                       + \frac{1}{8}
                       + \cdots
                       + \frac{1}{2^{n-1}}
                       + \frac{1}{2^n}} \\[2mm]
      \onslide<+->{& = \left(\frac{1}{2}
                     + \frac{1}{4}
                     + \frac{1}{8}
                     + \cdots
                     + \frac{1}{2^{n-1}}\right)
                     + \frac{1}{2^n}} \\[2mm]
    \onslide<+->{& = {\color<5->{blue}S_{n-1}} + \frac{1}{2^n}\\[2mm]}
    \onslide<+->{& = {\color<5->{blue}1 - \frac{1}{2^{n-1}}} + \frac{1}{2^n}\\[2mm]}
    \onslide<+->{& = 1 - \frac{2}{2^n} + \frac{1}{2^n}\\[2mm]}
    \onslide<+->{& = 1 - \frac{1}{2^n}}
  \end{align*}
\end{columns}
\end{frame}

%------------------------------------------------------------------------------%
\begin{frame}{Calculando o Limite das Somas Parciais}
    \begin{align*}
    	\onslide<+->{\lim_{n\to\infty} S_n}
      \onslide<+->{& = \lim \left(1 - \frac{1}{2^n}\right)     \\[2mm]}
    	\onslide<+->{& = \lim 1 - \lim\left(\frac{1}{2}\right)^{\!\!n} \\[2mm]}
    	\onslide<+->{& = 1 - 0 \\[2mm]}
      \onslide<+->{& = 1}
    \end{align*}

    \onslide<+->
    {
      Portanto a série converge e podemos escrever
      \[
        \sum_{n=1}^{\infty} \frac{1}{\;2^n\;} = 1
      \]
    }

\end{frame}

%------------------------------------------------------------------------------%
%------------------------------------------------------------------------------%
\section{Lista Mínima}

%------------------------------------------------------------------------------%
\begin{frame}{Lista Mínima}
  Estudar as Seção 6.1 da Apostila

  \vspace{\baselineskip}
  Exercícios: 1, 2, 4a-c, 6a-b

  \vfill
  Atenção: A prova é baseada no livro, não nas apresentações
\end{frame}

%------------------------------------------------------------------------------%
\end{document}
%------------------------------------------------------------------------------%
